\documentclass[12pt]{article}
\usepackage[margin=1in]{geometry}
\usepackage{amsmath, amssymb, amsthm, graphicx, hyperref}
\usepackage{subfigure}
\usepackage{color}
\usepackage{fancyhdr}
\usepackage{multirow, multicol}
\usepackage{pdfpages}
%\usepackage{stfloats}
\usepackage{tikz}
\pagestyle{fancy}

\fancypagestyle{firststyle}
{
  \renewcommand{\headrulewidth}{0pt}
  \fancyhead[LO]{\includegraphics{nyu_short_black}}
  \fancyhead[RO]{}
	\fancyfoot{}
}

\fancyhead[LO]{\scalebox{0.25}{\includegraphics{nyu_short_black}}}
\fancyhead[RO]{\tiny MATH-UA.120.001}


%\fancyfoot[RO]{\begin{tikzpicture}
%[scale=1, line/.style={black}]
%\draw[line] (0, 0)--(2,0)--(2, 2)--(0, 2)--(0, 0);
%%\node at (-0.5, 1) [anchor=east]{\large Score};
%\end{tikzpicture}}

\newcommand*\mycirc[1]{%
  \begin{tikzpicture}[baseline=(C.base)]
    \node[draw,circle,inner sep=1pt](C) {#1};
  \end{tikzpicture}}

\newcommand{\R}{\mathbb{R}}
\newcommand{\E}{\mathbb{E}}
\newcommand{\proj}{\textbf{proj}}
\newcommand{\curl}{\textrm{curl }}
\newcommand{\divergence}{\textrm{div }}
% For boldface vectors:
\renewcommand{\vec}[1]{\mathbf{#1}}
\renewcommand\arraystretch{2}

\newtheorem*{proposition}{Proposition}


\begin{document}
\thispagestyle{firststyle}
~
\vspace{0.5cm}

\begin{center}
\textbf{\Large MATH-UA.120.001 - Discrete Mathematics \\ \vspace{0.5cm} Final Exam 2016 Fall}\\~\\
Thursday, Dec 17, 2016
\end{center}

\vspace{1cm}

%\noindent
%\textbf{Name:} \underline{\hspace{15cm}}

\vspace{1cm}

\noindent
This exam is scheduled for 110 minutes, to be done individually, without calculators, notes, textbooks, and other outside materials.
\vspace{0.5cm}

\noindent
\textbf{Show all work to receive full credit, except where specified.}

\vspace{1cm}

\begin{multicols}{2}
\noindent

\noindent
\fbox{\parbox{10cm}{
\vspace{1cm}
\noindent
\textbf{Name: }

\vspace{1cm}

\noindent
\textbf{NYU NetID (email): }


\vspace{2cm}

\noindent
\emph{I pledge that I have observed the NYU honor code, and that I have neither given nor received unauthorized assistance during this examination.}
\noindent



\noindent

\vspace{1cm}

\noindent
\textbf{Signature:}

\vspace{2cm}
}}
%
%\vspace{0.5cm}
%

\begin{flushright}
\begin{tabular}{| c | c |}
\hline
Problem & Points \\
\hline
% & \\
TF & ~~~~~~~~~/20\\[0.5cm]
%& \\
\hline
% & \\
1 & ~~~~~~~~~/20\\[0.5cm]
%& \\
\hline
% & \\
2 & ~~~~~~~~~/15\\[0.5cm]
%& \\
\hline
% & \\
3 & ~~~~~~~~~/20\\[0.5cm]
%& \\
\hline
% & \\
4 & ~~~~~~~~~/25\\[0.5cm]
%& \\
\hline
\hline %& \\
Total & ~~~~~~~~~/100\\[0.5cm]
%& \\
\hline
\end{tabular}
\end{flushright}

\end{multicols}

\newpage

\section*{Answer Sheet (for Problem 1 True/False)} \label{answersheet}

%(3 points each) Select the correct option.  You need not explain your answer.  No partial credit will be given.
%
%\vspace{0.5cm}

\noindent
\textbf{You must record your final answers on the following answer sheet.  Only this page will be graded.}

\vspace{1cm}

\begin{center}
\begin{tabular}{| c | c c c c c |}
\hline

1 & \mycirc{T} & \mycirc{F} & & &  \\
% & & & & & \\
2 & \mycirc{T} & \mycirc{F} & & &   \\
% & & & & & \\
3 & \mycirc{T} & \mycirc{F} & & &   \\
% & & & & & \\
4 & \mycirc{T} & \mycirc{F} & & &  \\
% & & & & & \\
5 & \mycirc{T} & \mycirc{F} & & &  \\
6 & \mycirc{T} & \mycirc{F} & & & \\
7 & \mycirc{T} & \mycirc{F} & & & \\
8 & \mycirc{T} & \mycirc{F} & & & \\
9 & \mycirc{T} & \mycirc{F} & & & \\
10 & \mycirc{T} & \mycirc{F} & & & \\
\hline
\end{tabular}
\end{center}
\begin{flushright}
\begin{tikzpicture}
[scale=1, line/.style={black}]
\draw[line] (0, 0)--(2,0)--(2, 2)--(0, 2)--(0, 0);
\node at (-0.5, 1) [anchor=east]{\Large TOTAL};
\end{tikzpicture}
\end{flushright}
\renewcommand\arraystretch{1}

\begin{center}
\textbf{Make sure that you have marked your final answers on the answer sheet on page \pageref{answersheet}. Only answers on that page will be graded.}
\end{center}

\newpage
\section*{True or False}

(2 points each) Determine if each of the following statement is TRUE or FALSE.\\Choose TRUE if the statement is always true, and choose FALSE if it is not always true. Each question is worth 3 points.

\vspace{0.2cm}

\noindent
\textbf{Indicate your solution in the answer sheet on page \pageref{answersheet}.  You need not provide any justification.}


\begin{enumerate}
\item There are integers $x$ and $y$ satisfying the equation $15x+18y=9$. TRUE, $gcd(15,18)=3|9$
\vspace{0.8cm}
\item There is an element of $\mathbb{Z}_{12}$ that does not have a reciprocal. (Hint: Trying to compute the inverse of all elements is a very inefficient approach.) TRUE, e.g. $6$ as $gcd(6,12)=6\neq 1$.
\vspace{0.8cm}
\item Prime factorization of integers is easy and therefore it is easy to break public key encryption. FALSE
\vspace{0.8cm}
\item If $G$ is a subgraph of $H$, then $\alpha(G)\leq\alpha(H)$. FALSE, consider $G$ be an edgeless graph, while $H$ being the complete graph on the same vertices.
\vspace{0.8cm}
\item A group is a pair $(G,*)$ where $G$ is a set and $*$ is an operation on it that is commutative but not necessarily associative. FALSE, the operation in the group has to be associative.
\vspace{0.8cm}
\item No tree with more than one vertex has an Eulerian trail. FALSE, a path is a tree
\vspace{0.8cm}
\item Let $G$ be a graph. Then either $G$ or $\bar{G}$ must be connected. TRUE, it was a practice problem on HW13
\vspace{0.8cm}
\item If a graph has the same number of vertices as edges then it is a cycle.  FALSE, draw a triangle with an extra leg, for example
\vspace{0.8cm}
\item There is a graph with exactly one vertex with an odd degree. FALSE, it would make the sum of the degrees odd.
\vspace{0.8cm}
\item $\chi(G)\leq\Delta(G)$ for all graphs $G$.
\end{enumerate}
\begin{center}
\textbf{Make sure that you have marked your final answers on the answer sheet on page \pageref{answersheet}. Only answers on that page will be graded.}
\end{center}
\newpage
\section*{Free Response}

\textbf{Show all work and justification.} 

\begin{enumerate}
\item 
\begin{enumerate}
\item (3 points) Give the precise definition of the greatest common divisor of two positive numbers $a$ and $b$.\vspace{1cm}

The greatest common divisor of two positive integers $a$ and $b$ is an integer $c$ such that $c|a$ and $c|b$, and such that for every $d$ with the property that $d|a$ and $d|b$, we have $d\leq c$.

\item (7 points) Use Euclid's algorithm to compute $gcd(230,123)$. (No partial credit will be given for guessing.)
HELP: $107=6\cdot 16+11$
\begin{align*}
230 &= 1\cdot 123 + 107\\
123&= 1\cdot 107 + 16\\
107&= 6\cdot 16 + 11\\
16&= 1\cdot 11 + 5\\
11&= 2\cdot 5 +1\\
5&= 5\cdot 1 + 0
\end{align*}
and thus $gcd(230,123)=1$.

\vspace{0.5cm}
\item (7 points) Compute the reciprocal of $123$ in $\mathbb{Z}_{230}$. \vspace{0.5cm}
\begin{align*}
1 &= 11-2\cdot 5 = 11-2\cdot(16-11)=3\cdot 11 - 2\cdot 16=3\cdot(107-6\cdot 16)-2\cdot 16=\\
&3\cdot 107-20\cdot 16=3\cdot 107-20\cdot(123-107)=23\cdot 107-20\cdot 123=\\
&23\cdot(230-123)-20\cdot 123=23\cdot 230-43\cdot 123
\end{align*}
and therefore $123^{-1}=-43\mathrm{~mod~} 230 = 187$.
\item (3 points) Solve the equation $123\otimes x=2$ in $\mathbb{Z}_{230}$.
\[
x= 187\otimes 2 = 374\mathrm{~mod~}230=144
\]
\end{enumerate}\newpage
\newpage
\item In this problem we prove the following proposition.
\begin{proposition}
 Prove that consecutive perfect squares are relatively prime.
\end{proposition}

\begin{enumerate}
\item (5 points) First show that if there are two numbers $a$ and $b$ such that no prime divides them both, then $gcd(a,b)=1$.

If there are no prime that divides both of them then every prime in the prime factorizations has exponent zero at least in one of them. Therefore, in the prime factorization of $gcd(a,b)$, all primes are with a zero exponent and thus $gcd(a,b)=1$.
\item (8 points) Assume FTSC that there are consecutive perfect squares and a prime that divides them both. Reach a contradiction.
HINT: Can you drop the square in $p|a^2$?

Let the two consecutive perfect squares be $k^2$ and $(k+1)^2$. If $p|k^2$ then $p|k$ by the lemma from class. Similarly since $p|(k+1)^2$, we have $p|k+1$. These means that there are number $m$ and $n$ such that $k=mp$ and $k+1=np$. Substracting the two equations, we get $1=(n-m)p$ and thus $p|1$ contradicting $p$ being a prime.

\item (2 points) Conclude from this that the Proposition is true.

By part $b$, no primes divide both $k^2$ and $(k+1)^2$. By part a, this means that $gcd(k^2,(k+1)^2)=1$ which means that they are relative primes.
\end{enumerate}

\newpage
\item(20 points) Find all simultaneous solution of the following pair of congruences:
\begin{align*}
x&\equiv 5~({\rm mod}~9)\\
x&\equiv 7~({\rm mod}~16).
\end{align*}

By the first congruence,
\[
x = 9k +5
\]
where $k$ is an integer. By the second congruence,
\[
9k + 5 = x\equiv 7 (11)
\]
which can be also written as
\[
9k\equiv 2 (11).
\]
We find the solution in the form $k=k_0+11j$ for integer $j$ and $k_0$ being the solution of
\[
9\otimes k_0 = 2,\qquad {\rm in~\mathbf{Z}_11}.
\]
Since $9^{-1}=5$ in $\mathbf{Z}_11$, this means that
\[
k_0 = 5\otimes 2 = 10.
\]
Thus $k = 10 + 11j$ which we can plug back to the expression for $x$ to get
\[
x = 9(10+11j)+5=95+99j
\]

\newpage
\item \textbf{Graph theory}

Consider the following graph $G$:

\begin{figure}[h]
\centering
\includegraphics[scale=0.3]{Euler.png}
\end{figure}

\begin{enumerate}
\item (4 Points) Draw the subgraph $G[1,4,5,6]$ induced on the vertex set $\{1,2,4,5,6\}$.\vspace{3cm}
\item (4 Points) Find the minimum and maximum degree $\delta(G)$ and $\Delta(G)$.
\[
\delta(G)= 3,\qquad \Delta(G)=4
\]
\item (2 Points) Find a cut edge in $G[1,2,4,5,6]$.\vspace{1cm}
For example the edge ${2,4}$.
\item (4 Points) Does $G$ contain a spanning tree (with justification) and if yes find it (Draw. WARNING: Practice on the sketch paper, if there are multiple versions here, it will not be graded.)!\vspace{5cm}
Yes, the graph is connected. Break up cycles without loosing connectivity until you reach a tree.
\item (2 Points) Is this graph Eulerian? (with justification) If yes, exhibit an Euler tour (Draw).\vspace{1cm}
No, there are vertices with odd degrees.
\item (4 Points) Show that this graph is 4 colorable by exhibiting such a coloring. (Draw the graph again and indicate the colors by labels. WARNING: Practice on the sketch paper, if there are multiple versions here, it will not be graded.)\vspace{1cm} 
For example color $1,7$ as BLUE, $2,5$ as GREEN, $3,6$ as RED and $4$ as PURPLE.
\item (5 Points) Show that this graph is not $3$ colorable by assuming that it is and then arguing to a contradiction. HINT: Assign a color to vertex one and then discuss the colors of the other vertices in numerical order.

Assume the colors are RED, GREEN, and BLUE. Color the vertex $1$, let's say, BLUE. This forces $2$ or $3$ to be GREEN and RED or the other way around. In either case, $4$ must be colored BLUE. The same way $7$ must be colored BLUE. However, this necessarily creates a conflict along the edge $47$. This means that a proper coloring using three colors is not possible.

\end{enumerate}


\end{enumerate}

\newpage
\noindent
\textbf{Extra paper}
%
%\newpage
%\noindent
%\textbf{Extra paper}

%\includepdf[pages=1-2]{FormulaSheets.pdf}

\end{document}